\chapter{The Method of Least Squares}
\label{ch:leastsquares}

When two quantities are related—a wood's density and its hardness, calcium in
drinking water and a town's mortality—we often want to summarize the relationship
by a line and to use it for prediction. The method of \emph{least squares} fits
such a line by minimizing the total squared vertical distance between the data
and the line. Remarkably, this turns out to be maximum likelihood in disguise,
under the assumption that the measurement errors are normal.

\section{The Simple Linear Regression Model}

The \emph{simple linear regression model} for a bivariate dataset
$(x_1,y_1),\dots,(x_n,y_n)$ treats the $x_i$ as fixed and the responses as random:
\[
  Y_i=\alpha+\beta x_i+U_i,\qquad i=1,\dots,n,
\]
where $U_1,\dots,U_n$ are independent with mean $0$ and variance $\sigma^2$.
Equivalently, $Y_i\sim\N(\alpha+\beta x_i,\sigma^2)$: at each setting of $x$ there
is a whole population of possible responses, normally distributed about the line
$\E[Y]=\alpha+\beta x$, all sharing the same variance. The line's height shifts
with $x$, but the spread does not.

\section{Least Squares Is Maximum Likelihood}

To fit the line we choose $\alpha,\beta$ minimizing the sum of squared vertical
distances
\[
  S(\alpha,\beta)=\sum_{i=1}^n\bigl(y_i-\alpha-\beta x_i\bigr)^2.
\]
This is not an arbitrary choice. Under the normal model the loglikelihood is
\[
  \ell(\alpha,\beta,\sigma)=-n\ln\sigma-n\ln\sqrt{2\pi}-\frac{1}{2\sigma^2}\sum_{i=1}^n(y_i-\alpha-\beta x_i)^2,
\]
which, for any fixed $\sigma$, is maximized precisely when $S(\alpha,\beta)$ is
minimized. Least squares \emph{is} maximum likelihood for normal errors—one
reason normal-error models dominate applications. (Differentiating $\ell$ in
$\sigma$ also gives the MLE
$\hat\sigma=\sqrt{\tfrac1n\sum(Y_i-\hat\alpha-\hat\beta x_i)^2}$.)

As Figure~\ref{fig:ssparab} suggests, $S(\alpha,\beta)$ is an upward-opening
bowl (an elliptic paraboloid), so it has a unique minimum—except in the
degenerate case where all $x_i$ are equal.

\begin{figure}[ht]
  \centering
  \includegraphics[width=0.52\linewidth]{ss_paraboloid}
  \caption{The sum-of-squares surface $S(\alpha,\beta)$ is a bowl with a single
  lowest point $(\hat\alpha,\hat\beta)$.}
  \label{fig:ssparab}
\end{figure}

\section{The Normal Equations and the Estimates}

Setting the partial derivatives of $S$ to zero yields the \emph{normal
equations}
\[
  n\alpha+\beta\sum x_i=\sum y_i,\qquad
  \alpha\sum x_i+\beta\sum x_i^2=\sum x_iy_i,
\]
a $2\times 2$ linear system whose solution is the pair of least-squares estimates
\[
  \hat\beta=\frac{n\sum x_iy_i-(\sum x_i)(\sum y_i)}{n\sum x_i^2-(\sum x_i)^2},
  \qquad \hat\alpha=\ybar_n-\hat\beta\,\xbar_n.
\]
The second formula shows that the fitted line always passes through the center of
mass $(\xbar_n,\ybar_n)$ of the data.

\begin{example}[A three-point fit]
Fit the least-squares line to $(1,2),(3,1.8),(5,1)$, and find the residuals.

\solution With $\sum x_i=9$, $\sum y_i=4.8$, $\sum x_i^2=35$, $\sum x_iy_i=12.4$,
and $n=3$,
\[
  \hat\beta=\frac{3(12.4)-9(4.8)}{3(35)-9^2}=-0.25,\qquad
  \hat\alpha=1.6-(-0.25)(3)=2.35,
\]
so the line is $y=2.35-0.25x$ (Figure~\ref{fig:toyreg}). The residuals
$r_i=y_i-\hat\alpha-\hat\beta x_i$ are $-0.1,\,0.2,\,-0.1$, which sum to zero—as
they always must.
\end{example}

\begin{figure}[ht]
  \centering
  \includegraphics[width=0.42\linewidth]{toy_regression}
  \caption{The least-squares line for three points; the dotted segments are the
  residuals being minimized.}
  \label{fig:toyreg}
\end{figure}

\begin{example}[Predicting timber hardness]
For $36$ Australian hardwoods, regressing Janka hardness on density gives
$\hat\alpha=-1160.5$ and $\hat\beta=57.51$. Predict the hardness of a specimen of
density $65$.

\solution Substitute into the fitted line:
$\hat y=-1160.5+57.51(65)\approx 2577.7$. The regression line turns an
easily-measured density into a prediction of the hard-to-measure hardness.
\end{example}

\section{Unbiasedness of the Estimators}

The least-squares estimators are unbiased.

\begin{example}[$\hat\beta$ is unbiased]
Show $\E[\hat\beta]=\beta$ in the model $Y_i=\alpha+\beta x_i+U_i$.

\solution Since the $x_i$ are constants and $\E[Y_i]=\alpha+\beta x_i$, linearity
gives
\[
  \E[\hat\beta]=\frac{n\sum x_i\E[Y_i]-(\sum x_i)(\sum\E[Y_i])}{n\sum x_i^2-(\sum x_i)^2}
   =\frac{n\sum x_i(\alpha+\beta x_i)-(\sum x_i)\sum(\alpha+\beta x_i)}{n\sum x_i^2-(\sum x_i)^2}.
\]
The $\alpha$-terms cancel and the numerator collapses to
$\beta\bigl(n\sum x_i^2-(\sum x_i)^2\bigr)$, which is exactly the denominator
times $\beta$. Hence $\E[\hat\beta]=\beta$. It follows that
$\E[\hat\alpha]=\E[\Ybar_n]-\xbar_n\E[\hat\beta]=(\alpha+\beta\xbar_n)-\xbar_n\beta=\alpha$,
so $\hat\alpha$ is unbiased too.
\end{example}

For the error variance, the MLE $\tfrac1n\sum(Y_i-\hat\alpha-\hat\beta x_i)^2$
turns out to have expectation $\tfrac{n-2}{n}\sigma^2$, so the \emph{unbiased}
estimator divides by $n-2$ instead:
\[
  \hat\sigma^2=\frac{1}{n-2}\sum_{i=1}^n\bigl(Y_i-\hat\alpha-\hat\beta x_i\bigr)^2.
\]
We lose two degrees of freedom because two parameters, $\alpha$ and $\beta$, were
estimated from the same data.

\section{Residuals and Model Checking}

The \emph{residuals} $r_i=y_i-\hat\alpha-\hat\beta x_i$ are the leftover vertical
distances. They always sum to zero (because the line passes through
$(\xbar,\ybar)$), and plotting them against $x_i$ is the standard diagnostic. If
the linear model is right, the residual plot should look like featureless random
scatter about zero, as for the black cherry tree data
(Figure~\ref{fig:resgood}).

\begin{figure}[ht]
  \centering
  \includegraphics[width=0.55\linewidth]{residuals_good}
  \caption{Residuals with no trend: the linear model is adequate.}
  \label{fig:resgood}
\end{figure}

A \emph{pattern} in the residuals signals a misspecified model. For the timber
data, the residuals from a straight-line fit bend in a parabola
(Figure~\ref{fig:restimber}, left), suggesting a quadratic term:
\[
  Y_i=\alpha+\beta x_i+\gamma x_i^2+U_i.
\]
(This is still called a \emph{linear} model—linear in the parameters
$\alpha,\beta,\gamma$, even though the fitted curve is a parabola; see
Figure~\ref{fig:timberscatter}.) After adding the quadratic term, the residuals
lose their bend but begin to \emph{fan out} (Figure~\ref{fig:restimber}, right):
their spread grows with $x$. This is \emph{heteroscedasticity}—a violation of the
equal-variance assumption—as opposed to the constant-variance ideal of
\emph{homoscedasticity}. One remedy is weighted least squares, which we do not
pursue here.

\begin{figure}[ht]
  \centering
  \includegraphics[width=\linewidth]{residuals_timber}
  \caption{Left: a straight-line fit to the timber data leaves a parabolic
  trend. Right: after adding a quadratic term the trend is gone, but the
  residuals fan out, revealing heteroscedasticity.}
  \label{fig:restimber}
\end{figure}

\begin{figure}[ht]
  \centering
  \includegraphics[width=0.5\linewidth]{timber_scatter}
  \caption{Janka hardness versus density with a fitted quadratic. The model is
  linear in its parameters even though the curve is not a straight line.}
  \label{fig:timberscatter}
\end{figure}

\section{Regression Through the Origin}

Sometimes the intercept should be zero on physical grounds. For the black cherry
tree data, the response $y$ is timber volume and the predictor is $x=d^2h$
(diameter squared times height); as a tree shrinks to nothing, $x\to 0$ forces
$y\to 0$, so the model $Y_i=\beta x_i+U_i$ with no intercept is natural.

\begin{example}[The no-intercept slope]
Minimize $S(\beta)=\sum(y_i-\beta x_i)^2$ to derive the slope estimate.

\solution Differentiating, $S'(\beta)=-2\sum(y_i-\beta x_i)x_i=0$ gives
\[
  \hat\beta=\frac{\sum_{i=1}^n x_iy_i}{\sum_{i=1}^n x_i^2}. \qedhere
\]
\end{example}

\begin{example}[Three competing slope estimators]
For the no-intercept model, three natural estimators of $\beta$ are the average of
ratios $B_1=\tfrac1n\sum\tfrac{Y_i}{x_i}$, the ratio of averages
$B_2=\tfrac{\sum Y_i}{\sum x_i}$, and the least-squares estimator
$B_3=\tfrac{\sum x_iY_i}{\sum x_i^2}$. Show all three are unbiased, and compare
them on the cherry tree data.

\solution Each is a linear combination of the $Y_i$ with $\E[Y_i]=\beta x_i$, so
$\E[B_1]=\E[B_2]=\E[B_3]=\beta$: all unbiased. Numerically, using
$\sum x_i=87.456$, $\sum y_i=26.486$, $\sum y_i/x_i=9.369$, $\sum x_iy_i=95.498$,
$\sum x_i^2=314.644$, the three estimates are
\[
  B_1\approx 0.302,\qquad B_2\approx 0.303,\qquad B_3\approx 0.3035.
\]
They are close but not identical (Figure~\ref{fig:cherryreg}). Unbiasedness alone
does not single out a winner—efficiency does, and the least-squares estimator
$B_3$ is the one with smallest variance.
\end{example}

\begin{figure}[ht]
  \centering
  \includegraphics[width=0.5\linewidth]{cherry_regression}
  \caption{Black cherry tree volumes against $x=d^2h$ with the no-intercept
  least-squares line.}
  \label{fig:cherryreg}
\end{figure}

\section{Testing Whether a Coefficient Belongs}

A central question in regression is whether a coefficient is really nonzero, or
whether the apparent effect could be chance. We test, say, $H_0:\beta=\beta_0$
using a \emph{test statistic}
\[
  T_b=\frac{\hat\beta-\beta_0}{S_b},\qquad S_b^2=\frac{n}{n\sum x_i^2-(\sum x_i)^2}\,\hat\sigma^2,
\]
which standardizes the estimate by an estimate of its standard deviation. Under
$H_0$, $T_b$ has a $t$ distribution with $n-2$ degrees of freedom. Values near
zero support $H_0$; large values argue against it.

The strength of evidence is summarized by a \emph{$p$-value}—the probability,
\emph{if $H_0$ were true}, of a test statistic at least as extreme as the one
observed. A small $p$-value is strong evidence against $H_0$. If we reject when
the $p$-value falls below a threshold, that threshold is the \emph{significance
level}. A common informal scale:
\begin{center}
\renewcommand{\arraystretch}{1.15}
\begin{tabular}{@{}ll@{}}
\toprule
$p$-value & evidence against $H_0$\\
\midrule
$<0.01$ & very strong\\
$0.01$--$0.05$ & strong\\
$0.05$--$0.10$ & weak\\
$>0.10$ & little to none\\
\bottomrule
\end{tabular}
\end{center}

\begin{remark}
The $p$-value is \emph{not} $\Prob(H_0\text{ true}\mid\text{data})$—that
quantity belongs to Bayesian inference. The $p$-value is the smallest
significance level at which the data would lead us to reject $H_0$.
\end{remark}

\begin{example}[Does the cherry-tree intercept belong?]
Fitting the cherry tree data \emph{with} an intercept and testing $H_0:\alpha=0$
(we expect no intercept, since zero size should mean zero volume), the test
statistic is $t_a=-0.428$ on $29$ degrees of freedom, giving a one-sided
$p$-value $\Prob(t_{29}\le -0.428)=0.336$.

\solution This $p$-value is large, so we retain $H_0$: the data are entirely
consistent with a zero intercept, confirming the physical reasoning behind the
no-intercept model.
\end{example}

\begin{example}[Does calcium affect mortality?]
For $61$ English towns, mortality rate $y$ is regressed on calcium concentration
$x$ (Figure~\ref{fig:mortality}). We test $H_0:\beta=0$ against $H_1:\beta<0$
(more calcium, lower mortality). The fit gives $\hat\beta=-3.226$ and
$s_b=0.485$, so $t_b=-6.656$ on $59$ degrees of freedom.

\solution The corresponding $p$-value is about $5\times10^{-9}$—overwhelming
evidence against $H_0$. We conclude that higher calcium concentrations are
genuinely associated with lower mortality, not an artifact of chance.
\end{example}

\begin{figure}[ht]
  \centering
  \includegraphics[width=0.52\linewidth]{mortality_scatter}
  \caption{Mortality rate against calcium concentration for $61$ towns, with the
  least-squares line. The downward slope is highly significant.}
  \label{fig:mortality}
\end{figure}

\section{Reading a Regression Summary}

Statistical software reports several diagnostics for a fitted regression. The most
important is $R^2$, the \emph{coefficient of determination},
\[
  R^2=1-\frac{SS_{\mathrm{res}}}{SS_{\mathrm{tot}}},\qquad
  SS_{\mathrm{res}}=\sum_i(y_i-\hat y_i)^2,\quad SS_{\mathrm{tot}}=\sum_i(y_i-\ybar)^2,
\]
the fraction of the response's variation explained by the model. In simple linear
regression $R^2$ is just the square of the empirical correlation coefficient.
Figure~\ref{fig:rsq} illustrates the idea: $SS_{\mathrm{tot}}$ measures spread
about the flat mean line, $SS_{\mathrm{res}}$ about the fitted line; the more the
line shrinks the squares, the closer $R^2$ is to one.

\begin{figure}[ht]
  \centering
  \includegraphics[width=\linewidth]{r_squared}
  \caption{$R^2=1-SS_{\mathrm{res}}/SS_{\mathrm{tot}}$. The red squares (left)
  measure deviations from the mean; the blue squares (right) measure deviations
  from the fitted line. Smaller blue squares mean a better fit and larger $R^2$.}
  \label{fig:rsq}
\end{figure}

A summary also reports the \emph{adjusted} $R^2$, which penalizes extra
parameters and so guards against overfitting; the \emph{residual standard error},
$\hat\sigma$ on $n-2$ degrees of freedom (two parameters estimated); and an
\emph{$F$-test} of whether the regression coefficients are jointly zero. In
simple linear regression the $F$-test merely echoes the $t$-test for the slope,
but it becomes genuinely informative when several predictors are present.

Finally, $95\%$ confidence intervals for the parameters are obtained from the $t$
distribution with $n-2$ degrees of freedom. For the mortality data, the interval
for the calcium slope $\beta$ is about $(-0.0042,-0.0023)$—comfortably below
zero, consistent with the highly significant test above. For the quadratic timber
fit, the intercept's interval $(-804,556)$ comfortably contains zero, as physical
reasoning predicts, while the quadratic coefficient's interval
$(0.19,0.82)$ excludes zero, confirming that the curvature is real. Wide
intervals like the timber ones can be narrowed only by gathering more data.
